\subsection{Constructing the geometric filtration}
\label{sec:geometric-section-filtration}

The vector-space filtration in the determinant argument follows from
ordinary vanishing and the Koszul complex.  It does not require an
arithmetic extension estimate.

Let $Y$ be a smooth geometrically integral projective variety of
dimension $D\geq1$ over a number field $K$, and fix an integer
$T\geq D+1$.  Let $L$ be generated by
a fixed space of sections, let $A$ and $B$ be nef and big line bundles,
and fix
\[
 e\in H^0(Y,A\otimes B^{-1}\otimes L^{-D}),\qquad
 U\subset\{e\ne0\}
\]
with $U$ a nonempty open set.  Write $\mathcal G_j=\omega_Y\otimes A\otimes L^j$
and $G_j=H^0(Y,\mathcal G_j)$.  For $0\leq j\leq T$, suppose that the original spaces
$V_j\subset H^0(U,\mathcal G_j)$ are finite-dimensional and contain
$G_j$, that $V_0=G_0$, and
that multiplication by the fixed sections of $L$ maps $V_j$ into
$V_{j+1}$ for $j<T$.  Global sections are identified with their restrictions to
$U$, an injective map because $Y$ is integral.

\begin{proposition}\label{prop:geometric-section-filtration}
Let $s_1,\ldots,s_D$ in the fixed linear system
have a nonempty finite \'etale common zero scheme $Z$ contained in $U$.
Let $I$ be its ideal, set $a_0=T-D$, and put
\[
 P_s=\sum_{|\alpha|=a_0}s^\alpha V_D,\qquad
 W_s=G_T+P_s\subset V_T,\qquad
 F_a=W_s\cap H^0(U,I^a\mathcal G_T).
\]
After a finite extension splitting $Z$, with $\delta=\deg Z$, these
spaces and their natural symbol maps satisfy the filtration and product
identities in~\eqref{eq:arithmetic-generic-filtration}.  In particular,
\[
             F_{a_0}=P_s,
 \qquad F_a/F_{a+1}\simeq
           \bigoplus_{|\alpha|=a}E^\delta\quad(0\leq a<a_0).
\]
The evaluation map from $V_0$ onto the $\delta$ residue coordinates is
surjective.
\end{proposition}

\begin{proof}
Since $L$ is generated, it is nef.  For $C=A$ or $B$ and every $j\geq0$,
$C\otimes L^j$ is nef and big.  Kawamata--Viehweg vanishing therefore
gives
\[
       H^i(Y,\omega_Y\otimes C\otimes L^j)=0
                      \qquad(i>0,\ j\geq0).
\]
One may extend to an algebraic closure and descend this statement by
flat cohomology base change~\cite[Tag 02KH]{stacks-filtration-regular}.
We use precisely the line-bundle case of
\cite[Theorem 2.28]{viehweg-filtration-vanishing}.

At a point of $Z$, the $s_i$ are a regular system of parameters after
trivializing $L$.  Off $Z$ one is a unit.  Their Koszul complex is
therefore exact, and the associated graded algebra of $I$ is the
symmetric algebra on $I/I^2$; its degree-$a$ basis consists of the
monomials $s^\alpha$ with $|\alpha|=a$
\cite[Lemmas 31.21.1 and 31.21.4]{stacks-filtration-regular}.

First twist the Koszul complex by $\omega_Y\otimes B\otimes L^D$.
All its terms have the form $\omega_Y\otimes B\otimes L^{D-i}$,
where $0\leq i\leq D$, so the vanishing just proved makes restriction
onto $Z$ surjective.  Multiplication by $e$ maps these global sections
into $G_0=V_0$.  On $Z$, this multiplication is invertible.
Adjunction gives $(\omega_Y\otimes L^D)|_Z\simeq\omega_Z$, and
$\omega_Z\simeq\mathcal O_Z$ because $Z$ is \'etale.  Thus the actual
residue map $V_0\to H^0(Z,B|_Z)$ is onto.  No negative twist of $A$
has been used.

Next twist the Koszul complex by $\mathcal G_n$, with $n\geq D$.
Its terms have nonnegative exponents $n-i$, so it gives first-order
division:
\[
 H^0(Y,I\mathcal G_n)=\sum_{i=1}^D s_iG_{n-1}.
\]
Inductively, for every $a\geq0$ and $j\geq D$,
\begin{equation}\label{eq:geometric-global-division}
 H^0(Y,I^a\mathcal G_{j+a})
                  =\sum_{|\alpha|=a}s^\alpha G_j.
\end{equation}
Indeed write a section vanishing to order $a+1$ as
$\sum_{|\alpha|=a}s^\alpha g_\alpha$ with $g_\alpha\in G_{j+1}$,
using the induction hypothesis.  Its degree-$a$ symbol is zero.
Independence of the monomials in $I^a/I^{a+1}$ forces each
$g_\alpha|_Z=0$.  First-order division at $n=j+1$ now divides every
coefficient by one further $s_i$, with coefficients in $G_j$.
Collecting equal monomials proves the next step.  This proves the
required higher division without assuming vanishing for powers of $I$.

All elements of $P_s$ vanish to order $a_0$.  If $g+p\in W_s$ has
that order, with $g\in G_T$ and $p\in P_s$, then so does $g$.
Equation~\eqref{eq:geometric-global-division}, with $j=D$, puts $g$
in $\sum s^\alpha G_D\subset P_s$.  This proves $F_{a_0}=P_s$.

Finally split $Z$ and choose fibre coordinates for $B$ and $L$.
Write $L_p(l)$ for evaluation of a section $l$ at its $p$th point.
The local regular-sequence calculation identifies the degree-$a$
symbol of $s^\alpha l^{T-a}x$, for $x\in V_0$, with
\[
       (\mathbf1_{\alpha'=\alpha}
                       L_p(l)^{T-a}\operatorname{res}_p(x))_{\alpha',p}.
\]
Choose one $l$ over $K$ nonzero at every point of $Z$; this is possible
over the infinite field $K$ by avoiding finitely many hyperplanes.
The residue map is
onto, so these products supply every symbol.  They lie in $G_T$ and
hence in $W_s$.  The kernel of the symbol map is exactly $F_{a+1}$
by the definition of the $I$-adic filtration.  This supplies all the
claimed isomorphisms and identities.
\end{proof}

If the fixed linear system is birational onto its image $X$, good
sections as above exist on a nonempty open set of tuples.  Here
$\delta=\deg X$.  The projection construction used in
Proposition~\ref{prop:finite-section-jets} gives a morphism
$f:Y\to\mathbf P^D$ of generic degree $\delta$.  Let $U^\circ$ be
the intersection of $U$ with the birational isomorphism locus.
The image $f(Y\setminus U^\circ)$ is closed and has dimension at most
$D-1$.  Choose $q$ outside it and in the finite \'etale locus of $f$.
Its fibre consists of $\delta$ distinct points in $U^\circ$, and
its equations are $D$ sections from the same fixed linear system.

The universal common zero scheme over the affine tuple space is
proper.  At this tuple its relative Jacobian is invertible.  Remove
the closed images of its non-\'etale locus and of its intersection
with $Y\setminus U^\circ$.  Over the resulting nonempty open set,
the universal zero scheme is proper and \'etale, hence finite
\'etale, and has constant degree $\delta$.  Restrict to a nonempty
principal open.  All constructions are over $K$; since $K$ is
infinite, this open has $K$-rational tuples.

For the original Chow form $f$, its specialization at any such tuple
is nonzero and vanishes precisely on the hyperplanes through these
$\delta$ distinct image points.  Its degree is $\delta$, so
$f(s,l)=c_s\prod_pL_p(l)$ with $c_s\ne0$.  Thus the filtration and
Chow specialization required by the determinant theorem are supplied
by the original geometric hypotheses.  The arithmetic model, its
metrics, multiplication bounds, short integral families and geometric
height identity remain the data needed for the arithmetic application.
